\chapter{Transmission electron microscopy (TEM)}
\label{ch:tem}

\paragraph{In this chapter.} Electron--sample interactions, imaging modes (bright field, dark field/HAADF), particle-size
histograms, tomography, HRTEM and Fourier transforms (d-spacing $\to$ plane), analytical TEM (EDX, EELS, EFTEM) and
environmental TEM.

\section{Electron--beam/sample interactions}
\slide{189}: an incident electron beam produces back-scattered, secondary and Auger electrons, X-rays, visible light,
and transmitted electrons (forward, elastically and inelastically scattered).
\begin{classbox}
\begin{itemize}
  \item \textbf{Elastically} scattered electrons: ``conserve $E$, change direction'' $\to$ imaging, diffraction.
  \item \textbf{Inelastically} scattered: ``don't conserve $E$, change direction'' $\to$ EELS.
  \item Auger electrons: ``$e^-$ leaving by $E$ conservation when an $e^-$ downgrades to lower levels.''
  \item Secondary electrons: ``$e^-$ leaving by photons produced by $e^-$ downgrading to lower levels.''
\end{itemize}
\end{classbox}
Ruska: first electron microscope (Nobel 1986, shared with Binnig and Rohrer for STM) \slide{190}.
Instrument \slide{191}: electron gun, condenser lenses/aperture, specimen, objective lens/aperture, intermediate and
projector lenses, screen/camera; vacuum. Samples are dispersed (ultrasound in alcohol) on Cu grids with a lacy
carbon film \slide{192}.

\section{Imaging modes}
\slide{193}: the objective lens forms the \textbf{diffraction pattern} in its back focal plane and the \textbf{image}
in the image plane. An \textbf{objective aperture} in the back focal plane selects which electrons form the image.

\begin{keybox}[: bright field vs dark field]
\begin{tabularx}{\linewidth}{@{}lX@{}}
\toprule
\textbf{Bright field (BF)} \slide{194} & Aperture selects the \textbf{transmitted (direct) beam} (``you take $e^-$ in the
optical axis''). Heavy/thick regions scatter electrons out of the aperture $\Rightarrow$ appear \textbf{dark} on a
bright background. Example: Pt particles (dark dots) on CeO$_2$--ZrO$_2$ \slide{195}. \\
\textbf{Dark field (DF, STEM-HAADF)} \slide{196} & Aperture/annular detector collects \textbf{scattered} electrons
(``you take $e^-$ off the optical axis''); scattering regions appear \textbf{bright} on a dark background.
\textbf{HAADF} (high-angle annular dark field) intensity $\propto Z^2$: \textbf{Z-contrast}, heavy atoms bright.
Example: Au on TiO$_2$ (1--8\,\%), single Au atoms on CeO$_2$ \slide{197--200}. \\
\bottomrule
\end{tabularx}
\end{keybox}
\begin{tipbox}
HAADF/Z-contrast is the technique to \textbf{see heavy metal particles on a light support} (Pt, Au on oxides or carbon)
and to tell core from shell when the two metals have very different $Z$.
\end{tipbox}

\section{Particle size determination}
\begin{keybox}[: mean particle size \slide{201}]
Measure many particles (typically $>$100--200), build a histogram, and compute
\begin{equation}
\bar d = \frac{\sum_i n_i d_i}{\sum_i n_i},\qquad
\sigma = \sqrt{\frac{\sum_i n_i (d_i-\bar d)^2}{\sum_i n_i}}\ \text{\extra}
\end{equation}
$n_i$: number of particles with diameter $d_i$ (in a bin).
\end{keybox}
\begin{pitfall}
\slide{202}: a number-weighted mean hides rare big particles. One large particle can contain \textbf{50\,\% of the
metal mass} (w/w) while being invisible in the number histogram. Always report the \textbf{distribution} (mean $\pm$
deviation), and look at several areas of the grid. TEM sees few particles (not statistically representative), unlike
XRD or UV-vis.
\end{pitfall}

\section{TEM tomography}
\slide{203--204}: images recorded while tilting the sample (e.g.\ $-70^\circ$ to $+70^\circ$) are back-projected into a
\textbf{3D reconstruction}. It shows where particles are (inside pores of a zeolite, \textbf{inside or outside} a carbon
nanotube, PtRu/C), something a 2D projection cannot tell.

\section{High-resolution TEM (HRTEM) and Fourier transform}
\begin{keybox}[: HRTEM lattice fringes \slide{205--208}]
A large objective aperture lets the direct beam \emph{and} several diffracted beams interfere: the image shows
\textbf{lattice fringes}. ``Interference pattern, NOT atoms.'' The spacing between fringes $= d_{hkl}$ of the planes
parallel to the beam.\\
The \textbf{Fourier transform (FT/FFT)} of the image gives a diffraction-like pattern: each pair of spots corresponds to a
set of fringes. Spot distance $\propto 1/d$; angles between spots = angles between planes.\\
\textbf{Example} (Au along [110], \slide{208}): 2.04\,\AA{} $\to$ (200); 2.36\,\AA{} $\to$ (111) (two non-parallel sets).
\end{keybox}
\begin{recipe}{assigning HRTEM fringes}
\begin{enumerate}
  \item Measure the fringe spacing $d$ (or read it from the figure). Average over several fringes: $d = (\text{distance over }n\text{ fringes})/n$.
  \item Compute $d_{hkl} = a/\sqrt{h^2+k^2+l^2}$ for the \textbf{allowed} reflections of each candidate phase
        (fcc: 111, 200, 220, 311).
  \item Match within $\sim$2--3\,\%. Several phases? Pick the one whose $d$ matches. That identifies the particle.
  \item Angle between two fringe sets helps: e.g.\ two \{200\} sets at $90^\circ$ $\to$ viewed along [001];
        \{111\} and \{200\} at $54.7^\circ$; two \{111\} at $70.5^\circ$ (both along [110]).
  \item The fringes \textbf{parallel to a particle edge/surface} tell which plane is exposed at that surface.
\end{enumerate}
\end{recipe}
Particle shapes and exposed facets \slide{209}: cube = \{100\}; octahedron = \{111\}; intermediate truncated shapes
expose \{100\} + \{111\} (+\{110\}). Ceria cubes expose \{100\}, rods \{110\}+\{100\}, polyhedra \{111\}.
FT of SrTiO$_3$/TiO$_2$: 2.0\,\AA{} $\to$ SrTiO$_3$ (200) along [001]; 6.3/3.1/2.1\,\AA{} $\to$ Sr$_2$TiO$_4$ (002)/(004)/(006)
\slide{210}.

\subsection*{Epitaxy}
\begin{defbox}[: epitaxial relationship]
A particle grows in a defined crystallographic orientation with respect to its support: its lattice fringes are
\textbf{aligned/continuous} with those of the support across the interface (e.g.\ Cu\{hkl\} $\parallel$ CeO$_2$\{h'k'l'\}).
Epitaxy indicates a strong metal--support interaction (and usually small, flat, well-anchored particles). Large
particles that are randomly oriented with respect to the support are not epitaxial.\extra
\end{defbox}

\section{Analytical TEM: EDX, EELS, EFTEM}
\begin{itemize}
  \item \textbf{EDX (energy-dispersive X-ray analysis)} \slide{211}: the beam ejects a core electron; an outer electron
        fills the hole, emitting a characteristic X-ray (K$\alpha$, K$\beta$, L$\alpha$). Gives \textbf{elemental composition}
        of a selected area/particle and elemental \textbf{maps} (STEM-EDX). Class: ``At the same time you measure with
        $e^-$, you measure dispersed $e^-$ energy.'' Cu peaks often come from the grid.
  \item \textbf{EELS (electron energy loss spectroscopy)} \slide{212}: energy lost by transmitted electrons. Zero-loss peak
        (elastic), \textbf{plasmon} region (collective excitation of valence electrons, low loss $<$\SI{50}{eV}), and
        \textbf{core-loss} edges (inner-shell ionisation, element-specific, also sensitive to oxidation state).
        Low-loss spectra distinguish TiO$_2$ support from Au--Cu particles \slide{213}.
  \item \textbf{EFTEM (energy-filtered TEM)} \slide{214}: image formed only with electrons of a chosen energy loss
        $\Rightarrow$ \textbf{elemental maps} (e.g.\ Co red, O green: a Co particle with an oxide shell).
\end{itemize}
\begin{tipbox}
EDX/EELS maps answer ``which element is where'': core vs shell, alloy (both elements everywhere) vs segregation, and
the composition of a shell (Exercises~9b, 16).
\end{tipbox}

\section{Environmental TEM}
\slide{216--218}: gas (e.g.\ H$_2$) and heating inside the microscope; follows changes in real time: CeO$_2$ surface
reduction Ce$^{IV}$O$_2$ $\to$ Ce$^{III-IV}$O$_{2-x}$, Ni/Al$_2$O$_3$ growing a carbon nanotube in seconds \slide{217}.
Au nanoparticles grown on MWCNTs (functionalised with amines + AuCl$_4^-$ + citrate + ultrasound) \slide{215}.

\section{Advantages and limitations}
\proscons{
\begin{itemize}
  \item \textbf{Direct imaging}: size, shape, distribution, location of particles.
  \item Atomic-scale structure (HRTEM, FT): phases, exposed planes, epitaxy, defects.
  \item Local chemistry (EDX, EELS, EFTEM, HAADF), 3D (tomography), in situ (ETEM).
\end{itemize}}{
\begin{itemize}
  \item Local: few particles, not statistically representative.
  \item Vacuum, thin samples needed, beam damage (fullerenes fusing, \slide{121}).
  \item 2D projection: inside/outside ambiguous without tomography.
  \item Expensive; interpretation of HRTEM needs care.
\end{itemize}}
